\documentclass[11pt]{article}

\usepackage[margin=0.9in]{geometry}
\usepackage{amsmath,amssymb}
\usepackage{booktabs}
\usepackage{enumitem}
\usepackage{listings}
\usepackage{xcolor}
\usepackage{hyperref}

\definecolor{backcolour}{rgb}{0.95,0.95,0.92}
\definecolor{codeblue}{rgb}{0.05,0.15,0.55}
\definecolor{codegreen}{rgb}{0.0,0.45,0.0}

\lstset{
    backgroundcolor=\color{backcolour},
    basicstyle=\ttfamily\small,
    keywordstyle=\color{codeblue}\bfseries,
    commentstyle=\color{codegreen},
    stringstyle=\color{purple},
    breaklines=true,
    frame=single,
    showstringspaces=false,
    tabsize=4
}

\title{CPSC 250 \\ Operator Overloading}
\author{Edward Brash}
\date{}

\begin{document}

\maketitle

\section{Introduction}

Earlier, we discussed overloading comparison operators such as:

\[
<, >, ==, !=
\]

using methods such as:

\begin{lstlisting}[language=Python]
__lt__()
\end{lstlisting}

However, Python also allows us to overload arithmetic operators.

\section{Why Arithmetic Overloading Matters}

Suppose we define a \verb|Time| class.

We might naturally want to compute:

\begin{lstlisting}[language=Python]
time4 = time3 - time1
\end{lstlisting}

But what does subtraction mean for time objects?

Python does not know automatically.

Therefore, we must define the meaning ourselves.

\section{The \texttt{\_\_sub\_\_()} Method}

The subtraction operator:

\[
-
\]

is controlled by:

\begin{lstlisting}[language=Python]
def __sub__(self, other):
\end{lstlisting}

Here:

\begin{itemize}
    \item \verb|self| is the left object,
    \item \verb|other| is the right object.
\end{itemize}

\section{Example: Time Difference}

Suppose we define:

\begin{lstlisting}[language=Python]
class Time:

    def __init__(self, hours, minutes):
        self.hours = hours
        self.minutes = minutes
\end{lstlisting}

We can implement subtraction:

\begin{lstlisting}[language=Python]
def __sub__(self, other):

    time_diff = Time(0, 0)

    if self.minutes >= other.minutes:

        time_diff.minutes = self.minutes - other.minutes

        time_diff.hours = self.hours - other.hours

    else:

        time_diff.minutes = (
            self.minutes - other.minutes + 60
        )

        time_diff.hours = (
            self.hours - other.hours - 1
        )

    return time_diff
\end{lstlisting}

\section{Understanding the Logic}

Suppose:

\begin{lstlisting}[language=Python]
time1 = Time(10, 45)
time2 = Time(8, 20)
\end{lstlisting}

Then:

\begin{lstlisting}[language=Python]
time1 - time2
\end{lstlisting}

should represent:

\[
10:45 - 8:20 = 2:25
\]

If the minutes require borrowing, we borrow one hour and add 60 minutes.

\section{Returning a New Object}

Notice that the subtraction method creates a new \verb|Time| object:

\begin{lstlisting}[language=Python]
time_diff = Time(0, 0)
\end{lstlisting}

and returns it.

This is a common pattern in object-oriented programming.

\section{Other Numerical Operator Methods}

Python provides many special methods for arithmetic operations.

\begin{center}
\begin{tabular}{cc}
\toprule
Method & Operator \\
\midrule
\verb|__add__| & + \\
\verb|__sub__| & - \\
\verb|__mul__| & * \\
\verb|__truediv__| & / \\
\verb|__floordiv__| & // \\
\verb|__mod__| & \% \\
\verb|__pow__| & ** \\
\bottomrule
\end{tabular}
\end{center}

\section{Logical Operator Methods}

There are also methods associated with logical operations.

\begin{center}
\begin{tabular}{cc}
\toprule
Method & Meaning \\
\midrule
\verb|__and__| & logical and \\
\verb|__or__| & logical or \\
\bottomrule
\end{tabular}
\end{center}

\section{Conversion Methods}

Python also allows customization of built-in conversion functions.

\begin{center}
\begin{tabular}{cc}
\toprule
Method & Function \\
\midrule
\verb|__abs__| & \verb|abs()| \\
\verb|__int__| & \verb|int()| \\
\verb|__float__| & \verb|float()| \\
\bottomrule
\end{tabular}
\end{center}

\section{Example: Integer Conversion}

Suppose we define:

\begin{lstlisting}[language=Python]
class Time:

    def __int__(self):
        return 60*self.hours + self.minutes
\end{lstlisting}

Then:

\begin{lstlisting}[language=Python]
t = Time(2, 30)

print(int(t))
\end{lstlisting}

would produce:

\[
150
\]

because 2 hours and 30 minutes equals 150 total minutes.

\section{The Big Idea}

Operator overloading allows objects to behave naturally.

For example:

\begin{itemize}
    \item time subtraction behaves like subtraction,
    \item object comparison behaves like comparison,
    \item printing behaves naturally,
    \item conversions behave naturally.
\end{itemize}

This makes programs easier to read and easier to use.

\section{Key Takeaways}

\begin{itemize}
    \item Python allows arithmetic operator overloading.
    \item Special methods define the meaning of operators for custom classes.
    \item \verb|__sub__| defines subtraction.
    \item Many arithmetic operators can be overloaded.
    \item Built-in conversion functions can also be customized.
    \item Overloading helps objects behave naturally and intuitively.
\end{itemize}

\end{document}
